\scanpage{87}
induit l'identité sur $\varepsilon_a(D_0)$, il revient au
$=$ de se permettre de franchir les
positions supersingulières, or on a :
donc pour $n$ impair, $\pi_1(\widetilde{\mathcal X}\unsure{,|)}\longrightarrow\mathbb D_{2n}$
est aussi surjectif. Si $n$ est
pair, le résultat \inserted{sur $\varepsilon_a(D_0)$} d'un lacet $l$
\deleted{obtenu} ds $\mathcal X$ qui franchit
$\nu$ fois une position supersingulière,
est, \inserted{suivant que $\nu$ est pair ou impair,} le $=$ resp. opposé à celui obtenu
par le lacet $l'$ qui se déduit de
$l$ en remplaçant les franchissements
par des demi-tours. Si donc on
\deleted{prend} part d'une c. pos. rel. $D_0$ qui
passe par un sommet \inserted{$s_0$ de} $\Sigma$ \uline{d'ordre}
\uline{pair}, alors en modifiant le lacet $l$
demi-tour \inserted{de $D_0$} autour de $s_0$, de la
façon indiquée, on trouve une
\deleted{\unsure{autre}} action \deleted{\unsure{de}} $\varepsilon(D_0)$ qui n'est
pas l'identité. Cela prouve donc que
\uline{dans tous les cas}, $\pi_1(\widetilde{\mathcal X}\unsure{,|)}\longrightarrow\mathbb D_{2n}$ est
surjectif --- i.e. le revêt \inserted{principal associé} de $\unsure{\widetilde{\mathcal X}|\mathcal X}$, de groupe
$\mathbb D_{2n}$ est \uline{connexe}.
\margin{86}
